\documentclass[11pt]{article}
\usepackage{amsmath,amssymb}
\begin{document}
\section{Wolf Theorem 6.5 --- spectral radius as eigenvalue with PSD eigenvector}
\label{sec:gap_ch6_thm65}

\subsection{Source}
M.~Wolf, \emph{Quantum Channels \& Operations: Guided Tour}, Chapter~6,
Theorem~6.5, local source
\verb|Notes/WolfNoteTexSource/ch06_spectral_properties.tex|, lines~743--745.

\subsection{Statement}
Let $T:\mathcal{M}_d(\mathbb{C})\to\mathcal{M}_d(\mathbb{C})$ be a positive map
with spectral radius $\varrho$. Then $\varrho$ is an eigenvalue and there is a
positive semidefinite $X\in\mathcal{M}_d(\mathbb{C})$ such that $T(X)=\varrho X$.

\subsection{What is formalized}
\begin{itemize}
\item \texttt{exists\_posSemidef\_eigenvector\_general} (in
  \texttt{TNLean/Channel/PerronFrobenius/Existence.lean}) proves that every
  positive map has SOME nonnegative eigenvalue $r\ge 0$ with PSD eigenvector.
\item \texttt{spectralRadius\_eq\_of\_posDef\_eigenvector\_of\_irreducible\_cp}
  (in \texttt{TNLean/Channel/Irreducible/SpectralRadius.lean}) proves that for
  \emph{irreducible CP} maps, a positive eigenvalue with PosDef eigenvector equals
  the spectral radius. This is Wolf Theorem~6.3(4) but restricted to CP maps.
\end{itemize}

\subsection{What is missing}
\texttt{exists\_posSemidef\_eigenvector\_general} does not identify its eigenvalue
with the spectral radius (it may return a suboptimal eigenvalue, e.g.\ $r=0$ when
the map has a nontrivial kernel). The existing spectral-radius identity is
restricted to irreducible CP maps.

Proving the general positive-map case requires one of:
\begin{enumerate}
\item A direct resolvent/pole argument: for $\lambda>\varrho$, the resolvent
  $R(\lambda)=(\lambda-T)^{-1}$ preserves positivity; as $\lambda\downarrow\varrho$,
  the residue gives a PSD eigenvector for $\varrho$. Formalizing this needs
  holomorphic functional calculus and spectral projections.
\item Lifting from the irreducible case: decompose the PSD cone into invariant
  faces, apply the Perron--Frobenius theorem to each irreducible component, and
  select the component with maximal spectral radius. Requires formalizing the
  invariant-face decomposition of positive maps.
\end{enumerate}

\subsection{Elimination plan}
The extremal-quantity approach (issue \#3976 M3) will first prove Theorem~6.3
for irreducible positive maps directly (without CP/Kraus assumptions). Then the
invariant-face decomposition (a separate project) will lift to the general case.
If the resolvent argument is more tractable, it may be used as a shortcut.

\subsection{Verdict}
Wolf Theorem~6.5 is partially formalized: existence of a nonnegative PSD
eigenvalue is proved, but the identification with the spectral radius is missing.
\end{document}
